% ============================================================================
% Fichier  : partie0/P00_trois_jours_sans_banque.tex
% Titre    : Prologue -- Trois Jours sans Banque : quand Personne n'a Décidé
% Auteur   : MINKA MI NGUIDJOÏ Thierry Emmanuel
% Version  : 1.0 -- octobre 2026
% ============================================================================

\chapter[Prologue : Trois Jours sans Banque]{Prologue : Trois Jours sans Banque, quand Personne n'a Décidé}
\label{chap:prologue}
\chaptermark{Prologue}

\epigraph{%
    « Une panne se répare en quelques heures.\\
    Une décision que personne n'a prise se paie pendant des années. »%
}{Thierry Minka, auditeur des systèmes d'information}

\niveauChapitre{Introductif}{%
    Aucun prérequis. Ce prologue se lit avant d'ouvrir le moindre
    référentiel~: il ne suppose aucune connaissance de COBIT ni d'ITIL.%
}

\begin{boxobjectifs}
\begin{itemize}
    \item Reconstituer la chronologie d'une défaillance informatique majeure
          (opération SANAGA).
    \item Distinguer une cause technique d'une cause de gouvernance.
    \item Poser les trois questions réflexes~: qui décide, qui exécute, qui
          rend compte.
    \item Distinguer gouvernance et management des systèmes d'information.
    \item Identifier les parties prenantes qui demandent des comptes
          (conseil d'administration, régulateur, clients).
    \item Situer COBIT et ITIL dans la réponse à la défaillance.
\end{itemize}
\end{boxobjectifs}

\bigskip

% ============================================================================
\section*{Avant de commencer : une panne qui n'était pas technique}
% ============================================================================

Douala, lundi 9~mars~2026, 7~h~30. Les 38~agences de la BANQUE
SANAGA\footnote{La BANQUE SANAGA, ses dirigeants, ses prestataires, les
dates et les montants de ce cas sont fictifs. Les mécanismes de
gouvernance, les référentiels et les textes cités sont exacts.} ouvrent
leurs portes. Aucun guichet ne fonctionne. Les écrans affichent un message
d'erreur du nouveau système bancaire central, mis en service pendant le
week-end. Les distributeurs automatiques passent en mode dégradé, les
virements et les paiements par carte sont bloqués. À Yaoundé, à Bafoussam,
à Garoua, les files d'attente s'allongent devant des agences qui ne peuvent
ni verser ni retirer.

La paralysie dure \textbf{trois jours}. Le service ne revient complètement
que le jeudi 12~mars au matin. Le coût, chiffré plus tard par la direction
financière, atteint \textbf{1,9~milliard de francs CFA}. Le 16~mars, la
Commission bancaire de l'Afrique centrale (\sigle{COBAC}) demande à la banque
un rapport circonstancié sous trente jours.

\medskip

Le conseil d'administration réunit une séance extraordinaire. Le directeur
général, M.~FOTSO, ouvre la séance par une phrase qui résume ce que chacun
pense alors~: « Le prestataire nous a livré un logiciel défaillant. » Le
rapport d'enquête interne, remis trois semaines plus tard, dit tout autre
chose.

\begin{boxalert}
\textbf{Ce que révèle l'enquête.} Le logiciel n'était pas en cause. Il
avait été recetté à 63~\% seulement quand la date de bascule a été
maintenue. La date elle-même avait été fixée pour coïncider avec le
lancement commercial d'une offre mobile, déjà annoncée à la presse. Le
comité de pilotage du projet ne s'était plus réuni depuis quatre mois.
Aucun comité n'avait autorisé la mise en production. Aucun retour arrière
n'avait été testé, et la conversion des données de l'ancien système était
irréversible. Le contrat de l'intégrateur ne prévoyait aucune pénalité. Un
rapport d'audit interne avait signalé ces faiblesses cinq mois plus tôt~:
ses recommandations étaient restées « en cours ».

\smallskip
Chacune de ces défaillances est une \textbf{décision} qui n'a pas été
prise, ou qui a été prise par la mauvaise personne, sans trace écrite et
sans que personne en rende compte.
\end{boxalert}

\medskip

Ce prologue n'enseigne ni COBIT ni ITIL. Ces référentiels occupent les
chapitres qui suivent. Il enseigne \textbf{le réflexe fondateur de la
discipline}~: devant toute défaillance d'un système d'information, chercher
la décision derrière l'incident. Une panne se répare. Une organisation qui
ne sait pas qui décide reproduira la même panne.

\separateur

% ============================================================================
\section{Chronologie de la Défaillance}
\label{sec:pro-chronologie}
% ============================================================================

\subsection{Reconstituer avant de juger}

La première tâche de quiconque analyse une défaillance (le DSI qui rédige le
rapport au régulateur, l'auditeur mandaté par le conseil, le responsable des
risques) est de \textbf{reconstituer la chronologie}. Elle fixe les faits
avant toute interprétation. Elle se construit à partir de pièces~:
procès-verbaux, courriels, tickets, journaux d'exploitation, contrats.

\begin{boxdef}
\termecle{Chronologie de défaillance}\index{chronologie de défaillance}
(\textit{incident timeline})~: suite datée des faits, des décisions et des
absences de décision qui ont précédé, accompagné et suivi une défaillance.
Chaque entrée est rattachée à une pièce justificative.
\end{boxdef}

\begin{tableaucours}{L{2.6cm}L{6.4cm}L{3.2cm}}{%
    \tableauHeader{Date} &
    \tableauHeader{Fait établi} &
    \tableauHeader{Pièce}
}
Juin 2025 & Le conseil d'administration approuve le projet HORIZON
(remplacement du système bancaire central), budget de 3,2~milliards FCFA.
Intégrateur retenu~: DIGICORE AFRIQUE. & Procès-verbal du conseil \\
Oct. 2025 & L'audit interne signale l'absence de procédure de gestion des
changements et de plan de continuité testé. Recommandations notées « en
cours ». & Rapport d'audit n°~2025-07 \\
Nov. 2025 & Dernière réunion du comité de pilotage du projet. & Compte rendu \\
Janv. 2026 & Le directeur général fixe la bascule au week-end du 7--8~mars,
date du lancement de l'offre « SANAGA Mobile ». & Courriel du DG \\
20 févr. 2026 & Recette à 63~\% des scénarios. La DSI demande un report par
note interne. Aucune réponse écrite. & Note de la DSI \\
7 mars, 22 h & Début de la migration. & Journal d'exploitation \\
8 mars, 14 h & Conversion des comptes~: écarts de solde sur 11~\% des
comptes. Le retour à l'ancien système n'est plus possible. & Journal
d'exploitation \\
9 mars, 7 h 30 & Ouverture des agences~: système indisponible. & Tickets
du centre de services \\
12 mars, 8 h & Retour complet du service. & Journal d'exploitation \\
16 mars & La COBAC demande un rapport circonstancié sous 30~jours. &
Lettre de la COBAC \\
\end{tableaucours}
\vspace{-0.8em}
\begin{center}{\small\itshape Tableau~P.1. Chronologie de l'opération SANAGA}\end{center}

\subsection{Constatation et analyse}

Une chronologie ne contient que des \textbf{constatations}~: des faits
vérifiables, rattachés à une pièce. L'interprétation vient ensuite, dans une
partie distincte, et elle s'accompagne d'un niveau de confiance. Mélanger
les deux est l'erreur la plus fréquente des rapports de défaillance, et la
plus facile à exploiter par celui qui veut contester le rapport.

\begin{boxalert}
\textbf{Violation typique.} Le premier projet de rapport de la BANQUE
SANAGA écrit~: « Le prestataire a livré un logiciel non conforme, ce qui a
provoqué la panne. » La phrase mêle un fait non établi (la non-conformité)
et une causalité non démontrée. DIGICORE AFRIQUE produit aussitôt le
procès-verbal de recette, qui montre que la banque a elle-même décidé de
basculer avec 37~\% des scénarios non testés.

\smallskip
\textbf{Règle~:} écrire d'abord la constatation (« La recette couvrait
63~\% des scénarios le 20~février~; la bascule a eu lieu le 7~mars »), puis
l'analyse, avec son niveau de confiance (« Il est très probable que les
écarts de solde proviennent de scénarios de conversion non testés »).
\end{boxalert}

% ============================================================================
\section{Cause Technique, Cause de Gouvernance}
\label{sec:pro-causes}
% ============================================================================

\subsection{Deux niveaux d'explication}

Toute défaillance d'un système d'information s'explique à deux niveaux. Le
\textbf{niveau technique} répond à la question \emph{comment}~: quel
composant a cédé, quelle donnée a été corrompue. Le \textbf{niveau de la
gouvernance} répond à la question \emph{pourquoi cela a été possible}~:
quelle décision, quel contrôle, quelle responsabilité a manqué.

\begin{boxdef}
\termecle{Défaillance de gouvernance}\index{défaillance de gouvernance}
(\textit{governance failure})~: absence, inadéquation ou contournement d'un
mécanisme de décision, de contrôle ou de reddition de comptes, qui permet
qu'un risque connu ou connaissable se réalise.
\end{boxdef}

\begin{tableaucours}{L{5.6cm}L{6.6cm}}{%
    \tableauHeader{Cause technique (comment~?)} &
    \tableauHeader{Cause de gouvernance (pourquoi était-ce possible~?)}
}
Écarts de solde lors de la conversion des comptes & Bascule maintenue avec
une recette incomplète, sans décision formelle d'acceptation du risque \\
Retour arrière impossible & Aucun plan de retour arrière exigé ni testé
avant l'autorisation de mise en production \\
Indisponibilité prolongée & Aucun plan de continuité testé, malgré la
recommandation de l'audit interne \\
Intervention lente de l'intégrateur & Contrat sans engagement de niveau de
service ni pénalité \\
Communication de crise improvisée & Aucun rôle de crise défini, aucune
information préalable du conseil sur le risque du projet \\
\end{tableaucours}
\vspace{-0.8em}
\begin{center}{\small\itshape Tableau~P.2. Les deux niveaux d'explication de l'opération SANAGA}\end{center}

\subsection{Remonter de l'incident à la décision}

Une méthode simple permet de passer d'un niveau à l'autre~: demander
« pourquoi~? » jusqu'à atteindre une décision, ou une absence de décision.
Appliquée à la BANQUE SANAGA, elle donne~:

\begin{enumerate}
    \item \emph{Pourquoi les agences étaient-elles paralysées~?} Parce que
          le nouveau système produisait des soldes erronés.
    \item \emph{Pourquoi les soldes étaient-ils erronés~?} Parce que des
          scénarios de conversion n'avaient pas été testés.
    \item \emph{Pourquoi a-t-on basculé sans les tester~?} Parce que la date
          était fixée par le lancement commercial.
    \item \emph{Pourquoi la date l'a-t-elle emporté sur la recette~?} Parce
          que personne n'avait l'autorité, ni l'obligation, d'arbitrer entre
          le risque technique et l'enjeu commercial.
\end{enumerate}

La quatrième réponse n'est plus technique. Elle décrit une
\textbf{organisation}~: elle sera au cœur des Parties~\ref{part:fondements}
et~\ref{part:cobit} du cours.

\begin{boxinfo}
\textbf{La règle d'or~:} une analyse de défaillance est terminée quand elle
a atteint une décision, ou une absence de décision, et qu'elle a nommé le
rôle qui aurait dû la prendre. Tant qu'elle s'arrête à un composant, elle
n'a expliqué que la moitié de l'incident.
\end{boxinfo}

% ============================================================================
\section{Les Trois Questions Réflexes}
\label{sec:pro-trois-questions}
% ============================================================================

\subsection{Décider, exécuter, rendre compte}

Devant tout projet, tout changement ou toute défaillance d'un système
d'information, trois questions suffisent à révéler l'état de la gouvernance.

\begin{enumerate}
    \item \textbf{Qui décide~?} Quelle instance ou quelle personne avait
          l'autorité de trancher, et l'a-t-elle fait, par écrit~?
    \item \textbf{Qui exécute~?} Qui a mis en œuvre la décision, avec quels
          moyens, et selon quelle procédure~?
    \item \textbf{Qui rend compte~?} Qui doit expliquer le résultat, à qui,
          et sur la base de quelles informations~?
\end{enumerate}

\begin{boxdef}
\termecle{Droits de décision}\index{droits de décision} (\textit{decision
rights})~: répartition formelle, entre les instances et les personnes d'une
organisation, du pouvoir de décider sur une question donnée. Weill et Ross
en font le cœur de la gouvernance des SI~\cite{weill2004itgovernance}.

\medskip
\termecle{Redevabilité}\index{redevabilité} (\textit{accountability})~:
obligation, pour celui qui a reçu une responsabilité, d'en rendre compte à
l'instance qui la lui a confiée, et d'en assumer les conséquences.
\end{boxdef}

\subsection{Application à la BANQUE SANAGA}

\begin{tableaucours}{L{3.4cm}L{2.9cm}L{3.1cm}L{2.6cm}}{%
    \tableauHeader{Décision} &
    \tableauHeader{Prise par} &
    \tableauHeader{Aurait dû l'être par} &
    \tableauHeader{Trace écrite}
}
Date de bascule & Directeur général seul & Comité de pilotage, sur avis de
la DSI & Courriel \\
Bascule avec recette incomplète & Personne (statu quo) & Comité de
pilotage, avec acceptation formelle du risque & Aucune \\
Autorisation de mise en production & Personne & Comité des changements &
Aucune \\
Suite des recommandations d'audit & Personne & Direction générale, suivie
par le comité d'audit & « En cours » \\
Information du conseil sur le risque & Personne & Directeur général &
Aucune \\
\end{tableaucours}
\vspace{-0.8em}
\begin{center}{\small\itshape Tableau~P.3. Les droits de décision dans l'opération SANAGA}\end{center}

Le tableau parle de lui-même. Sur cinq décisions critiques, une seule a été
prise, par une personne seule, sans l'avis de ceux qui portaient le risque.
Les quatre autres n'ont été prises par personne. C'est la signature d'une
organisation sans gouvernance des SI~: non pas des dirigeants incompétents,
mais une absence de règles sur \emph{qui décide de quoi}.

\begin{boxmetier}{DSI}
\textbf{Protéger la décision, pas seulement le système}
\begin{enumerate}
    \item Tout désaccord sur un risque majeur (date, périmètre, recette) se
          formule \textbf{par écrit}, adressé à l'instance qui a le droit de
          décider, et non au seul supérieur hiérarchique.
    \item La note expose les faits (taux de recette, scénarios non testés),
          le risque (ce qui peut arriver, son coût estimé) et les options
          (reporter, réduire le périmètre, accepter le risque).
    \item Elle demande une \textbf{décision explicite}~: le silence n'est
          pas une décision.
    \item Le DSI conserve la note et la réponse, ou l'absence de réponse,
          dans le dossier du projet.
\end{enumerate}
La note du 20~février de la DSI de la BANQUE SANAGA était juste sur le fond.
Elle a échoué sur la forme~: adressée au seul directeur général, elle ne
demandait aucune décision et n'a jamais atteint le comité de pilotage.
\end{boxmetier}

% ============================================================================
\section{Gouvernance et Management}
\label{sec:pro-gouv-management}
% ============================================================================

\subsection{Une distinction fondatrice}

Les trois questions réflexes font apparaître deux familles d'activités. Les
unes fixent la direction et contrôlent le résultat~; les autres
réalisent. COBIT~2019 en fait une distinction structurante, chacune
relevant d'activités, de structures et de finalités
différentes~\cite{isaca2018cobitintro}.

\begin{boxdef}
\termecle{Gouvernance des SI}\index{gouvernance des SI} (\textit{governance
of enterprise I\&T})~: ensemble des activités par lesquelles l'organe de
gouvernance \emph{évalue} les besoins des parties prenantes pour fixer des
objectifs équilibrés, \emph{dirige} par la priorisation et la décision, et
\emph{surveille} la performance et la conformité au regard de cette
direction. Elle relève, dans la plupart des organisations, du conseil
d'administration.

\medskip
\termecle{Management des SI}\index{management des SI} (\textit{management
of enterprise I\&T})~: ensemble des activités qui \emph{planifient},
\emph{bâtissent}, \emph{exploitent} et \emph{surveillent} les systèmes
d'information, conformément à la direction fixée par la gouvernance. Il
relève de la direction générale.
\end{boxdef}

\begin{tableaucours}{L{3.0cm}L{4.6cm}L{4.6cm}}{%
    \tableauHeader{Critère} &
    \tableauHeader{Gouvernance} &
    \tableauHeader{Management}
}
Question & Fait-on les bonnes choses~? & Fait-on bien les choses~? \\
Responsable & Conseil d'administration et ses comités & Direction générale,
DSI, responsables opérationnels \\
Verbes (COBIT) & Évaluer, diriger, surveiller & Planifier, bâtir,
exploiter, surveiller \\
Horizon & Stratégique, pluriannuel & Opérationnel, annuel et quotidien \\
Exemple SANAGA & Décider si le lancement commercial justifie le risque de
bascule & Conduire la recette, préparer le retour arrière \\
\end{tableaucours}
\vspace{-0.8em}
\begin{center}{\small\itshape Tableau~P.4. Gouvernance et management des SI}\end{center}

\begin{boiteRemarque}{Le DSI gouverne-t-il~?}
La question revient à chaque promotion, et la réponse est nuancée. Le DSI
\emph{participe} à la gouvernance~: il éclaire le conseil, propose des
arbitrages, alerte sur les risques. Mais il n'en est pas l'organe. Un DSI
qui décide seul de la stratégie numérique gouverne par défaut, faute
d'instance compétente. Un conseil qui délègue tout au DSI renonce à sa
propre responsabilité. Le Chapitre~\ref{chap:gouv-entreprise} reprendra
cette question à partir de la théorie de l'agence.
\end{boiteRemarque}

% ============================================================================
\section{Ceux qui Demandent des Comptes}
\label{sec:pro-parties-prenantes}
% ============================================================================

\subsection{Les parties prenantes de la BANQUE SANAGA}

Une défaillance majeure fait apparaître, d'un coup, tous ceux à qui
l'organisation doit des comptes. Chacun pose une question différente et
attend une réponse documentée.

\begin{tableaucours}{L{3.6cm}L{4.4cm}L{4.2cm}}{%
    \tableauHeader{Partie prenante} &
    \tableauHeader{Sa question} &
    \tableauHeader{Ce qu'elle exige}
}
Conseil d'administration & Pourquoi n'avons-nous pas été informés du
risque~? & Un rapport sur les causes et la responsabilité de chacun \\
Actionnaires & Combien cela coûte-t-il, et cela peut-il se reproduire~? &
Le chiffrage des pertes et un plan d'action \\
COBAC & Le contrôle interne de la banque fonctionne-t-il~? & Un rapport
circonstancié sous 30~jours \\
Clients et commerçants & Nos opérations sont-elles justes~? & La
correction des soldes et une indemnisation \\
Commissaire aux comptes & Les comptes de l'exercice restent-ils
fiables~? & La preuve de l'exactitude des soldes convertis \\
\end{tableaucours}
\vspace{-0.8em}
\begin{center}{\small\itshape Tableau~P.5. Les parties prenantes et leurs exigences}\end{center}

\subsection{Le coût d'une décision manquée}

Le chiffrage de la direction financière donne la mesure de l'enjeu.

\begin{tableaucours}{p{8.2cm}r}{%
    \tableauHeader{Poste} &
    \tableauHeader{Montant (millions FCFA)}
}
Commissions et frais non perçus pendant trois jours & 410 \\
Indemnisation des clients et des commerçants & 380 \\
Prestations d'urgence et heures supplémentaires & 260 \\
Retraits hors ligne au-delà du solde, non recouvrés & 330 \\
Correction des données et audit externe imposé & 520 \\
\textbf{Total} & \textbf{1~900} \\
\end{tableaucours}
\vspace{-0.8em}
\begin{center}{\small\itshape Tableau~P.6. Coût de l'opération SANAGA}\end{center}

Le report de la bascule de six semaines, demandé par la DSI, aurait coûté
selon le même service environ 85~millions FCFA (prolongation de
l'intégrateur et décalage de la campagne commerciale). Le rapport entre les
deux montants, plus de vingt pour un, est l'argument que la gouvernance
aurait dû voir avant la bascule, et qu'elle n'a vu qu'après.

\subsection{Le cadre légal camerounais}

La défaillance de la BANQUE SANAGA ne relève pas seulement de la
réglementation bancaire. La loi camerounaise sur la cybersécurité soumet
les systèmes d'information à des obligations d'audit.

\begin{boxjuris}
\textbf{\loicam{n°~2010/012} du 21~décembre~2010 relative à la
cybersécurité et la cybercriminalité au Cameroun, article~32} (version
anglaise officielle)~:

\medskip
\textit{« (1) Electronic communication networks and information systems
shall be subject to a regime of compulsory and periodic auditing of their
security systems by the Agency. (2) Security audit and severity scale
rating shall be undertaken each year or as required by the prevailing
circumstances. »}

\medskip
Soit, en substance~: les réseaux de communications électroniques et les
systèmes d'information sont soumis à un audit obligatoire et périodique de
leur sécurité par l'Agence (l'\sigle{ANTIC}), conduit chaque année ou
selon les circonstances. L'article~13 de la même loi soumet en outre les
systèmes d'information des opérateurs à un audit de sécurité obligatoire,
dont les modalités sont fixées par voie réglementaire.

\medskip
\textbf{Portée.} La loi fait de l'audit des systèmes d'information une
obligation, et non une bonne pratique facultative. Le
Chapitre~\ref{chap:droit} précisera son champ d'application exact, ses
textes d'application et son articulation avec la réglementation de la
COBAC sur le contrôle interne des établissements de crédit, ainsi qu'avec
le droit OHADA de la responsabilité des dirigeants.
\end{boxjuris}

% ============================================================================
\section{Ce que COBIT et ITIL Auraient Changé}
\label{sec:pro-cobit-itil}
% ============================================================================

\subsection{Deux référentiels, deux niveaux}

Les deux référentiels au programme de ce cours répondent chacun à un niveau
d'explication de la défaillance. \textbf{COBIT~2019} (Partie~\ref{part:cobit})
organise la gouvernance et le management de l'ensemble des systèmes
d'information~: qui décide, selon quels objectifs, avec quels contrôles.
\textbf{ITIL~4} (Partie~\ref{part:pratique}) organise la gestion des
services informatiques au quotidien~: comment autoriser un changement,
traiter un incident, tenir un engagement de service. COBIT dit \emph{ce qui
doit être maîtrisé}~; ITIL décrit en détail \emph{comment le faire} pour
les services.

\begin{tableaucours}{L{4.2cm}L{4.2cm}L{3.8cm}}{%
    \tableauHeader{Défaillance SANAGA} &
    \tableauHeader{Objectif COBIT 2019} &
    \tableauHeader{Pratique ITIL 4}
}
Aucune instance n'arbitre le risque du projet & EDM01, EDM03 (cadre de
gouvernance, optimisation des risques) & · \\
Comité de pilotage en sommeil & BAI01, BAI11 (programmes, projets) &
Gestion de projet \\
Mise en production non autorisée, sans retour arrière & BAI06, BAI07
(changements, acceptation et transition) & Habilitation des changements,
gestion des mises en production \\
Contrat sans engagement de service & APO09, APO10 (accords de service,
fournisseurs) & Gestion des niveaux de service, gestion des
fournisseurs \\
Aucun plan de continuité testé & DSS04 (continuité) & Gestion de la
continuité de service \\
Recommandations d'audit sans suite & MEA02 (système de contrôle
interne) & · \\
\end{tableaucours}
\vspace{-0.8em}
\begin{center}{\small\itshape Tableau~P.7. Défaillances SANAGA, objectifs COBIT et
pratiques ITIL (intitulés traduits par nos soins)}\end{center}

\subsection{Un exemple~: la gestion des changements}

La mise en production du nouveau système était un \termecle{changement}
\index{changement} au sens des deux référentiels~: une modification d'un
élément du système d'information susceptible d'affecter les services.
COBIT en fait un objectif de management, que le cours présentera toujours
sous la même forme.

\begin{boxobjet}{BAI06 · Gérer les changements informatiques}
\textbf{Finalité~:} permettre des changements rapides et fiables, en
limitant le risque d'affecter la stabilité ou l'intégrité de
l'environnement modifié.

\textbf{Pratiques clés~:} évaluer, prioriser et autoriser les demandes de
changement~; gérer les changements d'urgence~; suivre l'état des
changements et en rendre compte~; clore et documenter les changements.

\textbf{Entrées et sorties~:} demande de changement, évaluation d'impact et
de risque, plan de retour arrière~; décision d'autorisation, calendrier des
changements, rapport de clôture.

\textbf{RACI (extrait)~:} autorisation d'un changement majeur~:
\emph{A} au comité des changements, \emph{R} au responsable du changement,
\emph{C} aux métiers concernés et au RSSI, \emph{I} au centre de services.

\textbf{Indicateurs~:} taux de changements en échec, nombre de changements
non autorisés, part des changements majeurs dotés d'un retour arrière
testé.
\end{boxobjet}

\begin{boxprocedure}
\textbf{Mise en production d'un changement majeur.} Pour qu'un auditeur
puisse établir que le changement était maîtrisé, le dossier doit contenir~:
\begin{itemize}
    \item le procès-verbal de recette, avec le taux de couverture des
          scénarios et la liste des anomalies ouvertes~;
    \item l'évaluation du risque et la décision écrite d'acceptation des
          risques résiduels, signée par l'instance compétente~;
    \item le plan de retour arrière et la preuve qu'il a été testé~;
    \item le procès-verbal d'autorisation du comité des changements, daté
          et antérieur à la mise en production~;
    \item le rapport de clôture, avec les incidents survenus.
\end{itemize}
Dans l'opération SANAGA, seule la première pièce existait.
\end{boxprocedure}

\begin{boxmetier}{Auditeur SI}
\textbf{Les cinq pièces à demander après une défaillance majeure}
\begin{enumerate}
    \item Les procès-verbaux des instances de gouvernance (conseil, comité
          d'audit, comité de pilotage) sur la période du projet.
    \item Le dossier du changement (voir l'exigence de contrôle ci-dessus).
    \item Les contrats avec les prestataires et leurs annexes de niveau de
          service.
    \item Le plan de continuité et les comptes rendus de ses derniers
          tests.
    \item Le suivi des recommandations d'audit antérieures, avec leurs
          dates d'échéance et leurs responsables.
\end{enumerate}
L'absence d'une pièce est elle-même une constatation~: elle se consigne
comme telle.
\end{boxmetier}

% ============================================================================
\section{Retour sur la BANQUE SANAGA : ce qui Aurait Dû se Passer}
\label{sec:pro-retour-sanaga}
% ============================================================================

Voici la même bascule, conduite cette fois par une organisation gouvernée.

\begin{boxscenario}{Opération SANAGA. La bascule gouvernée}
\begin{enumerate}
    \item \textbf{Juin 2025.} En approuvant le projet HORIZON, le conseil
          d'administration fixe ses règles~: un comité de pilotage mensuel
          présidé par le directeur général adjoint, un point de risque
          trimestriel au comité d'audit, et l'obligation d'une décision
          formelle avant toute bascule.
    \item \textbf{Octobre 2025.} Les recommandations de l'audit interne sur
          la gestion des changements et la continuité reçoivent chacune un
          responsable et une échéance. Le comité d'audit en suit l'avancement.
    \item \textbf{20 février 2026.} La note de la DSI (recette à 63~\%) est
          inscrite à l'ordre du jour du comité de pilotage. Le comité
          compare le coût d'un report (85~millions FCFA) et le risque d'une
          bascule prématurée.
    \item \textbf{Décision.} Le comité reporte la bascule de six semaines et
          dissocie le lancement de « SANAGA Mobile », qui démarre sur
          l'ancien système. La décision est consignée au procès-verbal.
    \item \textbf{Avril 2026.} Recette à 98~\%, retour arrière testé à
          blanc, comité des changements réuni. Le contrat de DIGICORE
          AFRIQUE a été amendé~: présence sur site pendant la bascule et
          pénalités de retard d'intervention.
    \item \textbf{Bascule.} Un écart de solde est détecté le dimanche à
          15~h. Le retour arrière est déclenché à 16~h. Les agences ouvrent
          le lundi sur l'ancien système. La bascule réussit deux semaines
          plus tard.
\end{enumerate}
\textbf{Résultat~:} le même incident technique s'est produit. Il est resté
un incident d'exploitation, sans conséquence pour les clients, au lieu de
devenir une crise à 1,9~milliard FCFA.
\end{boxscenario}

Le point essentiel est là~: \textbf{la gouvernance n'empêche pas les
incidents techniques}. Elle fait en sorte qu'ils soient anticipés, décidés
et contenus. La suite du cours donne les outils pour construire une telle
organisation.

% ============================================================================
\separateur
% ============================================================================

\begin{boxresume}
\textbf{Ce qu'il faut retenir de ce prologue~:}

\begin{itemize}
    \item \textbf{Une défaillance a deux niveaux}~: la cause technique dit
          comment le système a cédé, la cause de gouvernance dit pourquoi
          c'était possible. Une analyse qui s'arrête au composant est
          incomplète.
    \item \textbf{La chronologie d'abord}~: on reconstitue les faits à partir
          de pièces, et on sépare toujours la constatation de l'analyse.
    \item \textbf{Trois questions réflexes}~: qui décide, qui exécute, qui
          rend compte. Une décision sans trace écrite, ou prise par
          personne, est un signal d'alerte.
    \item \textbf{Gouvernance et management sont distincts}~: la gouvernance
          évalue, dirige et surveille~; le management planifie, bâtit,
          exploite et surveille.
    \item \textbf{Les parties prenantes demandent des comptes}~: conseil,
          actionnaires, régulateur, clients, commissaire aux comptes. Au
          Cameroun, la loi n°~2010/012 fait de l'audit des systèmes
          d'information une obligation.
    \item \textbf{COBIT et ITIL se complètent}~: COBIT organise ce qui doit
          être maîtrisé à l'échelle de l'organisation, ITIL décrit comment
          gérer les services au quotidien.
    \item \textbf{La gouvernance n'empêche pas l'incident}~: elle le rend
          anticipé, décidé et contenu.
\end{itemize}
\end{boxresume}

\bigskip

\begin{boxexo}[title={\ding{45}\quad Exercice~P.1. Technique ou gouvernance~? \niveaubase}]
Pour chacun des constats suivants, relevés dans d'autres organisations,
indiquez s'il s'agit d'une cause technique ou d'une cause de gouvernance,
et justifiez en une phrase.
\begin{enumerate}[(a)]
    \item Le disque d'un serveur de messagerie tombe en panne.
    \item Aucune sauvegarde de ce serveur n'avait été vérifiée depuis
          deux ans, sans que personne en soit chargé.
    \item Un administrateur applique un correctif un vendredi soir sans
          prévenir personne, conformément à l'usage du service.
    \item Le budget de renouvellement du matériel a été supprimé par le
          directeur financier sans consulter la DSI.
\end{enumerate}
\end{boxexo}

\begin{boxexo}[title={\ding{45}\quad Exercice~P.2. Les droits de décision de KRIBILOG \niveaumoyen}]
La société KRIBILOG SA, opérateur logistique à Kribi et à Douala, change
d'infogérant en trois semaines, sur décision de son directeur des achats,
pour une économie annuelle de 48~millions FCFA. Le responsable des services
informatiques l'apprend par le nouveau prestataire. Deux mois plus tard, le
logiciel de suivi des conteneurs connaît trois arrêts majeurs.
\begin{enumerate}[(a)]
    \item Construisez, sur le modèle du Tableau~P.3, le tableau des
          décisions critiques de ce changement de prestataire.
    \item Pour chaque décision, indiquez qui l'a prise, qui aurait dû la
          prendre, et quelle trace écrite devrait exister.
    \item Rédigez la constatation que l'auditeur SI consignerait, puis son
          analyse, avec un niveau de confiance.
\end{enumerate}
\end{boxexo}

\begin{boxexo}[title={\ding{45}\quad Exercice~P.3. Le rapport à la COBAC \niveauexpert}]
Vous êtes responsable GRC de la BANQUE SANAGA. Le directeur général vous
confie la rédaction du rapport circonstancié demandé par la COBAC.
\begin{enumerate}[(a)]
    \item Proposez le plan du rapport, en séparant la chronologie, les
          constatations, l'analyse des causes et le plan d'action.
    \item Le directeur général vous demande d'écrire que la panne est due
          au seul prestataire. Expliquez, pièces à l'appui, pourquoi cette
          affirmation ne résisterait pas à l'examen du régulateur.
    \item Pour chacune des six défaillances du Tableau~P.7, proposez une
          mesure corrective, son responsable et son échéance.
    \item Quelle place donnez-vous, dans ce rapport, à la recommandation de
          l'audit interne d'octobre~2025 restée sans suite~? Justifiez.
\end{enumerate}
\end{boxexo}

\bigskip

\begin{boxinfo}
\textbf{La suite du cours.} Ce prologue a posé le réflexe fondateur~:
chercher la décision derrière l'incident. La Partie~\ref{part:fondements}
donne à ce réflexe ses fondements, en commençant par la gouvernance
d'entreprise et la théorie de l'agence (Chapitre~\ref{chap:gouv-entreprise}).
La Partie~\ref{part:cobit} construit avec COBIT~2019 le système de
gouvernance qui a manqué à la BANQUE SANAGA, et la
Partie~\ref{part:pratique} en organise, avec ITIL~4, la pratique
quotidienne.
\end{boxinfo}
